\section{Introduction}
\label{sec:intro}

Vision-language-action (VLA) models map camera images and a language instruction directly to robot actions. Starting from robot transformers trained on large demonstration sets~\citep{rt1} and from fine-tuning a vision-language model to output actions as tokens~\citep{rt2}, the field has produced open generalist policies~\citep{octo,openvla}, faster fine-tuning recipes~\citep{oft} and flow-matching action experts~\citep{pi0,pi05,smolvla}. These models are compared mostly in simulation, and LIBERO~\citep{libero} has become the standard benchmark. OpenVLA-OFT alone reports an average of \ExtOftAverage\% over its four suites~\citep{oft}, so new methods are separated from old ones by a few points, and an audit of LIBERO results found that only \ExtJiangSigPct\% of claimed state-of-the-art improvements are provably statistically significant~\citep{jiang2026bench}.

A success rate is produced by an evaluator, not by the model alone. Before the policy sees a frame, the evaluator processes it with image steps copied from how the training data were built. In the official OpenVLA evaluation on LIBERO the simulator renders the camera at $\RenderRes \times \RenderRes$; step one encodes the frame as JPEG and decodes it again, as the training dataset builder did; step two resizes it to $\InputRes \times \InputRes$ with an antialiased Lanczos3 filter, following the Octo data loader~\citep{octo} that OpenVLA used for training; and step three takes a center crop of \CropAreaPct\% of the area, which the evaluation code requires for checkpoints trained with random-crop augmentation~\citep{openvlacode} (Figure~\ref{fig:pipeline}). Every new evaluation harness re-implements these steps, and a step that is dropped or changed raises no error: it only gives a different number. The vla-eval harness found that the undocumented crop step is worth about \ExtVeCropPts{} points of OpenVLA's LIBERO score~\citep{vlaeval}, while its reproduction notes change the other steps only together~\citep{vlaevalrepo}.

In our own reproduction of the published baselines, runs under MuJoCo \MujocoNewVersion{} instead of \MujocoVersion{} gave OpenVLA-7B \MujocoNewK{} instead of \MujocoOldK{} successful episodes of \PilotPairs{} on initial states \SplitA{} of LIBERO-Spatial, mostly in one task that fell from \MujocoOldTaskFive{} to \MujocoNewTaskFive{} of \PilotStates, and OpenVLA-OFT fell from \OftOldTaskFive{} to \OftNewTaskFive{} in the same task (Appendix~\ref{app:mujoco}). When the evaluation can move this much, two questions follow for anyone who compares checkpoints across harnesses: which image steps must be reproduced exactly, and how many episode outcomes any small change of the input reshuffles. Robustness benchmarks perturb the scene; here the scene is clean and only the evaluator's own processing changes.

We therefore change one evaluator image step at a time and measure the effect on closed-loop success of two released OpenVLA-family checkpoints, OpenVLA-7B~\citep{openvla} and OpenVLA-OFT~\citep{oft}, on all \NPairs{} initial states of LIBERO-Spatial. Every episode of a variant is paired with the reference episode that starts from the same initial state; an identical rerun measures the noise floor of our stack, and sham cells, meaningless changes as large as the mildest steps fixed in a public plan before they ran, measure how many flips such a change produces. For OpenVLA-7B the sham cells flip about as many episodes as the mild steps, so the main finding is the fragility of single outcomes, not the effect of a step; OpenVLA-OFT flips far less, and there the steps lose episodes the shams do not. Our contributions are:
\begin{itemize}
\item a paired, one-step-at-a-time ablation of the official OpenVLA evaluation image pipeline for two released OpenVLA-family checkpoints, on all \NPairs{} initial states of LIBERO-Spatial, with the primary comparison, its threshold and the decision rule written down before the full variant cells ran, after a pilot on initial states \SplitA{} that included two of the variant steps, and without a public timestamp; the \OutPairs{} initial states outside the pilot are reported separately (Section~\ref{sec:stats}, Appendix~\ref{app:pilot});
\item bounds rather than effects: no single step has a net effect that \NPairs{} pairs resolve (for JPEG off on OpenVLA-7B the one-sided \AlphaPct\% bootstrap bound on the loss is \SevJpegUB{} points), \HookContribFlips{} (Section~\ref{sec:results});
\item an exploratory open-loop check on identical observations, which separates what a step does to the predicted action from what the closed loop does with it (Section~\ref{sec:openloop});
\item per-episode logs, evaluation code and analysis scripts that regenerate every number in the paper.\footnote{Per-episode logs, evaluation harness, analysis scripts and paper source: \RepoURL.\HarnessNote}
\end{itemize}
The observation that opposing outcome changes cancel in an aggregate success rate, and the paired statistics we use to count them, are not new~\citep{liberoctrl,li2026paraphrase}, and image classification counts the same flips between two model versions as negative flips and prediction churn~\citep{yan2021pct,bhojanapalli2021churn}; \HookIntroNew
